%% Chapter 7: Student Account Management
\chapter{Student Account Management}
\label{ch:student-mgmt}

\section{Creating Student Accounts}

ALAMS stores all student accounts in the cloud database. Each student needs one account to log in from any lab.

\subsection{Method A: Admin Dashboard (Recommended for Bulk Setup)}

\begin{enumerate}[leftmargin=*, label=\textbf{Step \arabic*.}]
  \item Log in to the Admin Dashboard at \texttt{http://SERVER\_IP:3000/admin/login}.
  \item Click \textbf{Student Directory} in the left sidebar.
  \item Scroll to the \textbf{Add New Student} form.
  \item Fill in all fields:
    \begin{tcolorbox}[alamsbox, title={Student Registration Form Fields}]
    \begin{tabular}{ll}
    \textbf{Field} & \textbf{Value}\\
    Enrollment Number & Institution-assigned ID (e.g. \texttt{ENR2026001})\\
    Full Name & Student's full legal name\\
    Password & Minimum 8 characters (shared with mobile login)\\
    PIN & 6-digit numeric PIN (for offline/PIN fallback login)\\
    \end{tabular}
    \end{tcolorbox}
  \item Click \textbf{Add Student}.
  \item The student account is active immediately.
\end{enumerate}

\subsection{Method B: API (Bulk Import via Script)}

\begin{tcolorbox}[codebox, title={Bulk Student Creation via API}]
\begin{lstlisting}[language=bash]
# Create one student account
curl -X POST http://SERVER_IP:5000/api/v1/auth/signup \
  -H "Content-Type: application/json" \
  -d '{
    "enrollmentNumber": "ENR2026011",
    "fullName": "Riya Kapoor",
    "password": "Student@2026!",
    "pin": "123456",
    "role": "STUDENT"
  }'

# Expected: {"token":"...","user":{"id":"...","enrollmentNumber":"ENR2026011"}}
\end{lstlisting}
\end{tcolorbox}

\subsection{Method C: Modify the Seed File}

For a fresh deployment, add student objects to \texttt{server/prisma/seed.ts} in the \texttt{studentData} array and re-run the seed:

\begin{tcolorbox}[codebox, title={Add Students to Seed Data}]
\begin{lstlisting}[language=TypeScript]
const studentData = [
  // Existing students...
  { enrollmentNumber: "ENR2026011", fullName: "Riya Kapoor" },
  { enrollmentNumber: "ENR2026012", fullName: "Amit Kumar" },
  // Add as many as needed
];
\end{lstlisting}
\end{tcolorbox}

\section{Default Student Credentials}

All seeded pilot student accounts share these credentials:

\begin{tcolorbox}[successbox, title={\faKey\quad Pilot Student Credentials}]
\begin{center}
\begin{tabular}{rl}
\textbf{Enrollment Numbers:} & \texttt{ENR2026001} through \texttt{ENR2026010}\\
\textbf{Password (Mobile Login):} & \texttt{Student@2026!}\\
\textbf{PIN (Lock Screen Fallback):} & \texttt{123456}\\
\end{tabular}
\end{center}
\end{tcolorbox}

\begin{tcolorbox}[dangerbox, title={\faExclamationCircle\quad Before Real Deployment}]
\begin{itemize}[label={\faTimesCircle}]
  \item Change all student PINs to unique values per student.
  \item Change all passwords to institution-generated secure passwords.
  \item Communicate credentials to students individually (not publicly displayed).
\end{itemize}
\end{tcolorbox}

\section{Managing Student Status}

\subsection{Disable a Student Account}

If a student is suspended or should not have lab access:

\begin{enumerate}[leftmargin=*, label=\textbf{\arabic*.}]
  \item Dashboard → Student Directory
  \item Find the student row
  \item Toggle the \textbf{Status} switch to \textcolor{aurxonred}{\textbf{Inactive}}
  \item The student cannot authenticate (any method) until reactivated.
\end{enumerate}

\subsection{Student Password Reset}

Currently, password reset is done by creating a new account or using the API:

\begin{tcolorbox}[codebox, title={Simulate Password Reset (API)}]
\begin{lstlisting}[language=bash]
# With admin authorization, create a new account with updated password
# (Note: full reset endpoint is planned for v2.0)
curl -X POST http://SERVER_IP:5000/api/v1/auth/signup \
  -H "Content-Type: application/json" \
  -d '{"enrollmentNumber":"ENR2026001","fullName":"Arjun Sharma","password":"NewPass@2026!","pin":"654321","role":"STUDENT"}'
\end{lstlisting}
\end{tcolorbox}

\section{Student Enrollment Number Standards}

It is recommended to adopt a consistent naming convention for enrollment numbers:

\begin{table}[h!]
\centering
\caption{Recommended Enrollment Number Format}
\label{tab:enrollment-format}
\begin{tabularx}{\textwidth}{lX}
\toprule
\textbf{Format} & \textbf{Description}\\
\midrule
\texttt{ENR\{YEAR\}\{3-digit-seq\}} & E.g. \texttt{ENR2026001}, \texttt{ENR2026099}\\
\rowcolor{aurxonlightgray}
\texttt{ADMIN\{2-digit\}} & E.g. \texttt{ADMIN01}, \texttt{ADMIN02}\\
\texttt{SUPER\{2-digit\}} & E.g. \texttt{SUPER01}\\
\rowcolor{aurxonlightgray}
\texttt{FAC\{3-digit\}} & E.g. \texttt{FAC001}, \texttt{FAC002}\\
\bottomrule
\end{tabularx}
\end{table}
